\chapter{群论的应用}\label{ch:b3:group-theory-applied}

在空气中，每一百万个分子里二氧化碳只占几百个，
其余几乎全是氮气和氧气；可吸收温暖地面射向太空的
红外线的，却是二氧化碳，而不是氮气或氧气。
分子只有通过使偶极矩发生变化的振动，才能吸收红外线，
而振动是否具有这一性质，
完全由对称性决定：氮气的唯一振动不能，
二氧化碳的四种振动中有两种能。
本章把\cref{ch:b3:point-groups}中的\omterm{def:b3:point-groups:irrep}{特征标表}化为工具：
用来约化表示、构造对称性匹配轨道、统计并标记振动，
并导出红外光谱与拉曼光谱的选择定则。

\begin{recall}
\Cref{ch:b3:point-groups}定义了\omterm{def:b3:point-groups:operation}{对称操作}、\omterm{def:b3:point-groups:point-group}{点群}、
类、表示、\omterm{def:b3:point-groups:representation}{特征标}和\omterm{def:b3:point-groups:irrep}{特征标表}。本科二年级分册
构造了 \ce{H2O}、\ce{NH3} 和 \ce{CH4} 的分子轨道，
利用氢 $1s$ 轨道的片段轨道和对称性匹配组合，
并以本章推导的标记为之命名。本科一年级分册
以波数读取红外光谱，并定义了分子的
极化率。
\end{recall}

\section{约化表示}

附着于分子的任意一组函数或矢量，都携带其\omterm{def:b3:point-groups:point-group}{点群}的一个表示。
其中大多数为\omterm{def:b3:group-theory-applied:reducible}{可约表示}。

\begin{definition}[可约表示、全对称表示]\label{def:b3:group-theory-applied:reducible}
称\emph{可约表示}\index{可约表示}为其全部矩阵都能经一次基变换
化为相同分块对角形式的表示，
此时它是若干\omterm{def:b3:point-groups:irrep}{不可约表示}之和，
记作 $\Gamma = \sum_in_i\Gamma_i$。其\omterm{def:b3:point-groups:representation}{特征标}全为 1 的\omterm{def:b3:point-groups:irrep}{不可约表示}
即\emph{全对称表示}\index{全对称表示}，
记为（$A_1$、$A_{1g}$、
$A_1'$\dots）。
\end{definition}

\begin{theorem}[大正交定理]\label{thm:b3:group-theory-applied:great-orthogonality}
设 $\Gamma_i$、$\Gamma_j$ 为两个\omterm{def:b3:point-groups:irrep}{不可约表示}，维数分别为 $d_i$、
$d_j$，群的阶为 $h$，并用酉矩阵写出，
\[
\sum_R D_i(R)_{mn}^*\,D_j(R)_{m'n'} = \frac{h}{d_i}\,\delta_{ij}\,\delta_{mm'}\,\delta_{nn'}.
\]
\end{theorem}
\admitted

证明借助舒尔引理，将在更高级的课程中讨论；其
对\omterm{def:b3:point-groups:representation}{特征标}的推论正是化学所用的内容，本书每一张\omterm{def:b3:point-groups:irrep}{特征标表}
都依据这些推论加以检验。

\begin{corollary}[特征标的正交性]\label{cor:b3:group-theory-applied:characters}
记 $g_c$ 为类 $c$ 中操作的个数，
\[
\sum_cg_c\,\chi_i(c)^*\chi_j(c) = h\,\delta_{ij},
\]
则\omterm{def:b3:point-groups:irrep}{不可约表示}的个数等于类的个数，
且 $\sum_id_i^2 = h$。
\end{corollary}

\begin{proof}
在定理中取 $m = n$ 与 $m' = n'$，并对 $m$、$m'$ 求和：左边化为
$\sum_R\chi_i(R)^*\chi_j(R)$，右边为 $\frac{h}{d_i}\delta_{ij}
\sum_{m}\delta_{mm} = h\delta_{ij}$；按类对操作分组。表的各行
因而是维数等于类数的空间中的正交向量，
故\omterm{def:b3:point-groups:irrep}{不可约表示}的个数至多等于类数；
等号成立及 $\sum d_i^2 = h$，由将定理用于正则表示得到
（见练习 10）。
\end{proof}

\begin{theorem}[约化公式]\label{thm:b3:group-theory-applied:reduction}
设表示 $\Gamma$ 的\omterm{def:b3:point-groups:representation}{特征标}为 $\chi(c)$，它包含\omterm{def:b3:point-groups:irrep}{不可约表示}
$\Gamma_i$ 恰好
\[
n_i = \frac1h\sum_cg_c\,\chi(c)\,\chi_i(c)^*
\]
次。这就是\emph{约化公式}\index{约化公式}。
\end{theorem}

\begin{proof}
由 $\Gamma = \sum_jn_j\Gamma_j$，其\omterm{def:b3:point-groups:representation}{特征标}为 $\chi(c) = \sum_jn_j\chi_j(c)$。
两边乘以 $g_c\chi_i(c)^*$，对类求和，并利用上述推论：
$\sum_cg_c\chi(c)\chi_i(c)^* = \sum_jn_jh\delta_{ij} = hn_i$.
\end{proof}

\begin{example}[氨的氢原子轨道]\label{ex:b3:group-theory-applied:nh3}
取 \ce{NH3} 中三个氢原子的 $1s$ 轨道作为基。操作 $E$
保持三个轨道都不动（$\chi = 3$），$C_3$ 使三个轨道都移位（$\chi = 0$），
每个 $\sigma_v$ 保持一个不动、交换另两个（$\chi = 1$）。取 $C_{3v}$ \omterm{def:b3:point-groups:irrep}{特征标表}：
$n_{A_1} = \frac16(3 + 0 + 3\times1) = 1$, $n_{A_2} = \frac16(3 + 0 - 3) = 0$,
$n_E = \frac16(6 + 0 + 0) = 1$。于是 $\Gamma_{\mathrm H} = A_1 + E$：三个氢原子
给出一个全对称组合和一个简并对。
\end{example}

\begin{method}[约化表示]\label{met:b3:group-theory-applied:reduce}
\begin{enumerate}
  \item 对每一类求\omterm{def:b3:point-groups:representation}{特征标}：对一组以原子为中心的函数，
        即保留在原位的函数个数，
        每个函数按其获得的符号计入。
  \item 对每个\omterm{def:b3:point-groups:irrep}{不可约表示}，用其\omterm{def:b3:point-groups:representation}{特征标}行应用\omterm{thm:b3:group-theory-applied:reduction}{约化公式}，
        以各类的大小为权重。
  \item 检验：结果为非负整数，且 $\sum_in_id_i$
        等于 $E$ 下的\omterm{def:b3:point-groups:representation}{特征标}（即基的大小）。
\end{enumerate}
\end{method}

\section{对称性匹配轨道}

\begin{definition}[投影算符]\label{def:b3:group-theory-applied:projection}
称映到\omterm{def:b3:point-groups:irrep}{不可约表示} $\Gamma_i$ 上的\emph{投影算符}\index{投影算符}
为
\[
\hat P_i = \frac{d_i}{h}\sum_R\chi_i(R)^*\,\hat R,
\]
其中 $\hat R$ 将操作 $R$ 作用于一个函数。
\end{definition}

\begin{proposition}[投影算符的作用]\label{prop:b3:group-theory-applied:projection}
对任意函数 $f$，$\hat P_if$ 或为零，或为一个按
$\Gamma_i$ 变换的函数（对称性匹配组合）；对一维的 $\Gamma_i$，
它至多相差符号 $\chi_i(S)$ 而保持不变，
对每个操作 $S$ 成立。
\end{proposition}

\begin{proof}[部分证明]
当 $d_i = 1$ 时：$\hat S\hat P_if = \frac1h\sum_R\chi_i(R)\hat S\hat Rf$。令 $R' = SR$，
它与 $R$ 一样遍历整个群；$\chi_i(R) = \chi_i(S^{-1}R') =
\chi_i(S)\chi_i(R')$，对一维表示（此时\omterm{def:b3:point-groups:representation}{特征标}即为
矩阵本身），于是 $\hat S\hat P_if = \chi_i(S)\hat P_if$。一般情形
需用大正交定理，在本阶段不加证明地承认。
\end{proof}

\begin{method}[构造对称性匹配组合]\label{met:b3:group-theory-applied:salc}
\begin{enumerate}
  \item 任取该组中的一个函数 $h_1$，
        列出每个操作将其映为的函数。
  \item 对每个\omterm{def:b3:point-groups:irrep}{不可约表示}，将各像乘以该操作的
        \omterm{def:b3:point-groups:representation}{特征标}，求和并归一化（略去原子函数之间的
        重叠）。
  \item 对简并表示，再对另一函数作投影，以得到
        第二个成员，然后正交化。
\end{enumerate}
\end{method}

\begin{example}[氨与水的组合]\label{ex:b3:group-theory-applied:salcs}
在 \ce{NH3} 中，$\hat P_{A_1}h_1 \propto h_1 + h_2 + h_3$（\omterm{def:b3:point-groups:representation}{特征标}均为 1）。对 $E$
（\omterm{def:b3:point-groups:representation}{特征标}为 $2, -1, 0$），$\hat P_Eh_1 \propto 2h_1 - h_2 - h_3$；对 $h_2$ 作投影并
正交化后得到伙伴 $h_2 - h_3$。归一化后：$a_1 =
\frac{1}{\sqrt3}(h_1 + h_2 + h_3)$，$e = \frac{1}{\sqrt6}(2h_1 - h_2 - h_3)$，
$\frac{1}{\sqrt2}(h_2 - h_3)$。在 \ce{H2O} 中，将分子置于 $xz$ 平面，$h_1 + h_2$ 为
$a_1$，$h_1 - h_2$ 为 $b_1$（在 $C_2$ 及 $\sigma_v'(yz)$ 作用下变号，
这两个操作交换两个氢原子）。若按某些书的习惯，将分子置于 $yz$ 平面，
则全文中 $b_1$ 与 $b_2$ 的标记互换。
\end{example}

\begin{omfigure}
\begin{tikzpicture}[font=\small]
  \foreach \x/\lab/\ca/\cb/\cc in {0/{$a_1$}/1/1/1, 4/{$e$}/2/-1/-1, 8/{$e$}/0/1/-1}{
    \begin{scope}[xshift=\x cm]
      \foreach \a/\c in {90/\ca,210/\cb,330/\cc}{
        \pgfmathsetmacro{\r}{0.18+0.13*abs(\c)}
        \ifdim \c pt > 0pt \fill[omDef!55, draw=omDef] (\a:1.0) circle (\r);
        \else \ifdim \c pt < 0pt \fill[omThm!25, draw=omThm] (\a:1.0) circle (\r);
        \else \draw[gray, densely dotted] (\a:1.0) circle (0.12); \fi\fi
      }
      \node[aN, minimum size=4.6mm] at (0,0) {};
      \node[anchor=north] at (0,-1.45) {\lab};
    \end{scope}}
  \node[anchor=north, font=\footnotesize] at (0,-1.85) {$h_1 + h_2 + h_3$};
  \node[anchor=north, font=\footnotesize] at (4,-1.85) {$2h_1 - h_2 - h_3$};
  \node[anchor=north, font=\footnotesize] at (8,-1.85) {$h_2 - h_3$};
\end{tikzpicture}

{\small 氨分子中三个氢原子 $1s$ 轨道的对称性匹配组合，
沿 $C_3$ 轴俯视（$h_1$ 在顶部）。蓝色：系数为正，
红色：系数为负，圆盘大小随
系数绝对值增大；虚线：系数为零。}
\end{omfigure}

\begin{example}[金属周围的六个配体]\label{ex:b3:group-theory-applied:ml6}
对从八面体顶点指向金属的六个配体 $\sigma$ 轨道，
其在 $O_h$ 群的各类操作（$E$，$8C_3$，$6C_2$，$6C_4$，
$3C_2$，$i$，$6S_4$，$8S_6$，$3\sigma_h$，$6\sigma_d$）上的\omterm{def:b3:point-groups:representation}{特征标}依次为 $6, 0, 0, 2, 2, 0, 0, 0, 4,
2$，约化得 $\Gamma_\sigma = A_{1g} + E_g + T_{1u}$。金属的 $s$
轨道（$a_{1g}$）、$d_{z^2}$ 与 $d_{x^2-y^2}$（$e_g$）及 $p$（$t_{1u}$）轨道均能找到配体
伙伴；其 $d_{xy}$、$d_{xz}$、$d_{yz}$（$t_{2g}$）轨道则无伙伴，
在 $\sigma$ 中保持非键——这正是本科二年级分册中 $t_{2g}$/$e_g$ 分裂的来源，
并在\cref{ch:b3:organometallic-bonding,ch:b3:complex-spectra-magnetism}中展开。
\end{example}

\begin{omfigure}
\begin{tikzpicture}[font=\small, yscale=0.95]
  % O orbitals (left), MOs (centre), H combinations (right); schematic energies
  \foreach \y/\l in {0/{$2s\ (a_1)$}}{\draw[omstate] (0,\y) -- (1.2,\y); \node[anchor=east] at (-0.05,\y) {\l};}
  \draw[omstate] (0,3.0) -- (0.36,3.0) (0.42,3.0) -- (0.78,3.0) (0.84,3.0) -- (1.2,3.0);
  \node[anchor=east] at (-0.05,3.0) {$2p\ (b_1, a_1, b_2)$};
  \draw[omstate] (6.8,3.9) -- (8.0,3.9); \node[anchor=west] at (8.05,3.9) {$h_1 - h_2\ (b_1)$};
  \draw[omstate] (6.8,3.6) -- (8.0,3.6); \node[anchor=west] at (8.05,3.5) {$h_1 + h_2\ (a_1)$};
  \foreach \y/\l in {-0.6/{$2a_1$}, 1.6/{$1b_1$}, 2.4/{$3a_1$}, 3.0/{$1b_2$}, 5.6/{$4a_1$}, 6.2/{$2b_1$}}{
    \draw[omstate] (3.4,\y) -- (4.6,\y); \node[anchor=west] at (4.65,\y) {\l};}
  \foreach \y in {-0.6,1.6,2.4,3.0}{\node[font=\scriptsize] at (4.0,\y+0.16) {$\uparrow\downarrow$};}
  \draw[densely dotted, gray] (1.2,0) -- (3.4,-0.6) (1.2,3.0) -- (3.4,1.6) (1.2,3.0) -- (3.4,2.4) (1.2,3.0) -- (3.4,3.0)
     (6.8,3.6) -- (4.6,-0.6) (6.8,3.9) -- (4.6,1.6) (6.8,3.6) -- (4.6,2.4) (1.2,3.0) -- (3.4,6.2) (6.8,3.9) -- (4.6,6.2) (6.8,3.6) -- (4.6,5.6);
  \node[anchor=north] at (0.6,-1.0) {O};
  \node[anchor=north] at (4.0,-1.0) {\ce{H2O}};
  \node[anchor=north] at (7.4,-1.0) {2 H};
  \draw[->] (-2.3,-0.9) -- (-2.3,6.4) node[above] {$E$};
\end{tikzpicture}

{\small 水的价层分子轨道示意图（示意能量，分子位于
$xz$ 平面内）。只有对称性相同的轨道才会混合：$b_1$ 型氢组合
与 $2p_x$ 混合，$a_1$ 型组合则与 $2s$ 及 $2p_z$ 混合；$2p_y$（$b_2$）无伙伴，
保持为非键孤电子对，是最高占据轨道。}
\end{omfigure}

\section{振动模式}

\begin{definition}[简正模式]\label{def:b3:group-theory-applied:normal-mode}
分子的\emph{简正模式}\index{简正模式}是一种集体
振动，其中所有原子以相同频率同相振动，
沿质量加权黑塞矩阵的特征向量振动
（见\cref{prop:b3:computational-chemistry:hessian}）。每个简正模式
按\omterm{def:b3:point-groups:point-group}{点群}的一个\omterm{def:b3:point-groups:irrep}{不可约表示}变换。
\end{definition}

\begin{proposition}[全部运动的表示]\label{prop:b3:group-theory-applied:gamma-3n}
在每个原子上附加三个位移矢量 $x$、$y$、$z$。对某一操作
$R$，这一 $3N$ 维表示的\omterm{def:b3:point-groups:representation}{特征标}为
\[
\chi_{3N}(R) = (\text{保持不动的原子数}) \times \chi_{xyz}(R),
\]
其中 $\chi_{xyz} = 1 + 2\cos\theta$ 对应旋转角 $\theta$，$-1 + 2\cos\theta$ 对应
非真旋转（$E$ 为 3，$-1$ 对应 $C_2$，0 对应 $C_3$，1 对应 $\sigma$，
$-3$ 对应 $i$）。
\end{proposition}

\begin{proof}
移到别处的原子，其矢量移出了矩阵的对角线，
对迹没有贡献。留在原位的原子，
其三个矢量经 $3\times3$ 矩阵变换，该矩阵由 $R$ 给出，其迹
在取 $z$ 沿轴的坐标系中计算：旋转时为 $\cos\theta + \cos\theta + 1$，
非真旋转时末项变为 $-1$。
\end{proof}

\begin{proposition}[振动计数]\label{prop:b3:group-theory-applied:mode-count}
$\Gamma_{\mathrm{vib}} = \Gamma_{3N} - \Gamma_{\mathrm{trans}} - \Gamma_{\mathrm{rot}}$，其中
$\Gamma_{\mathrm{trans}}$ 为 $(x, y, z)$ 的表示，$\Gamma_{\mathrm{rot}}$ 为
$(R_x, R_y, R_z)$ 的表示；非线性分子有 $3N - 6$ 个振动，线性
分子则为 $3N - 5$ 个。
\end{proposition}

\begin{proof}
这 $3N$ 个位移描述了全部运动：质心的三个
平动、三个转动（线性分子为两个，因绕自身轴转动
不使任何原子核移动），以及振动。这些子空间在各操作下
互不混合，故其表示相加。
\end{proof}

\begin{example}[水]\label{ex:b3:group-theory-applied:water}
% ledger: vib:H2O.sym, vib:H2O.bend, vib:H2O.asym
\ce{H2O} 属 $C_{2v}$，分子位于 $xz$ 平面内。留在原位的原子数为 3、1、3、1；
$\chi_{xyz}$：$3, -1, 1, 1$；故 $\chi_{3N} = 9, -1, 3, 1$。约化得 $3A_1 +
A_2 + 3B_1 + 2B_2$。平动 $A_1 + B_1 + B_2$ 与转动 $A_2 + B_1 + B_2$
扣除后，$\Gamma_{\mathrm{vib}} = 2A_1 + B_1$。两个 $a_1$ 模式分别是
对称伸缩（\qty{3657}{cm^{-1}}）与弯曲（\qty{1595}{cm^{-1}}）；$b_1$ 模式是
反对称伸缩（\qty{3756}{cm^{-1}}）。
\end{example}

\begin{omfigure}
\begin{tikzpicture}[font=\small, >=Stealth]
  \foreach \x/\name/\lab in {0/sym/{$\nu_1$, $a_1$, 对称伸缩}, 4.3/bend/{$\nu_2$, $a_1$, 弯曲}, 8.6/asym/{$\nu_3$, $b_1$, 反对称伸缩}}{
    \begin{scope}[xshift=\x cm]
      \node[aO] (o\name) at (0,0.35) {};
      \node[aH] (a\name) at (-0.8,-0.25) {};
      \node[aH] (b\name) at (0.8,-0.25) {};
      \begin{scope}[on background layer]\draw[bond] (o\name.center) -- (a\name.center) (o\name.center) -- (b\name.center);\end{scope}
      \node[anchor=north, align=center, font=\footnotesize] at (0,-0.9) {\lab};
    \end{scope}}
  % symmetric stretch: H outward along the bonds, O up a little
  \draw[->, omThm, thick] (-0.95,-0.36) -- (-1.42,-0.71);
  \draw[->, omThm, thick] (0.95,-0.36) -- (1.42,-0.71);
  \draw[->, omThm, thick] (0,0.74) -- (0,0.98);
  % bend: each H moves perpendicular to its bond, closing the angle; O up
  \draw[->, omThm, thick] (4.3-0.92,-0.36) -- (4.3-0.62,-0.76);
  \draw[->, omThm, thick] (4.3+0.92,-0.36) -- (4.3+0.62,-0.76);
  \draw[->, omThm, thick] (4.3,0.74) -- (4.3,0.98);
  % antisymmetric stretch: one H out, the other in, O sideways
  \draw[->, omThm, thick] (8.6-0.95,-0.36) -- (8.6-1.42,-0.71);
  \draw[->, omThm, thick] (8.6+1.02,0.0) -- (8.6+0.70,0.24);
  \draw[->, omThm, thick] (8.6-0.05,0.80) -- (8.6+0.25,0.80);
\end{tikzpicture}

{\small 水的三种\omterm{def:b3:group-theory-applied:normal-mode}{简正模式}（箭头表示位移，未按比例绘制）。
两个 $a_1$ 模式保持分子的全部对称性；$b_1$
模式在 $C_2$ 作用下改变符号。}
\end{omfigure}

\begin{example}[另外四种分子]\label{ex:b3:group-theory-applied:table}
同样的步骤（在本章的图形数据中计算并核对）
给出：
\begin{center}\small
\begin{tabular}{llll}
\hline
分子 & \omterm{def:b3:point-groups:point-group}{点群} & $\Gamma_{\mathrm{vib}}$ & 模式数 \\ \hline
\ce{NH3} & $C_{3v}$ & $2A_1 + 2E$ & 6 \\
\ce{CH4} & $T_d$ & $A_1 + E + 2T_2$ & 9 \\
\ce{CO2} & $D_{\infty h}$（经 $D_{2h}$）& $\Sigma_g^+ + \Sigma_u^+ + \Pi_u$（$A_g + B_{1u} + B_{2u} + B_{3u}$）& 4 \\
\ce{BF3} & $D_{3h}$ & $A_1' + 2E' + A_2''$ & 6 \\
\ce{XeF4} & $D_{4h}$ & $A_{1g} + B_{1g} + B_{2g} + A_{2u} + B_{2u} + 2E_u$ & 9 \\
\ce{SF6} & $O_h$ & $A_{1g} + E_g + T_{2g} + 2T_{1u} + T_{2u}$ & 15 \\ \hline
\end{tabular}
\end{center}
对线性分子，无限群应换为其\omterm{def:b3:point-groups:class}{子群} $D_{2h}$，
其中简并的 $\Pi_u$ 弯曲模式分裂为 $B_{2u} + B_{3u}$。
\end{example}

\section{选择定则}

\begin{definition}[直积]\label{def:b3:group-theory-applied:direct-product}
两个表示的\emph{直积}\index{直积} $\Gamma_i\otimes\Gamma_j$ 是
由它们基函数之积所承载的表示，
其\omterm{def:b3:point-groups:representation}{特征标}为 $\chi_i(R)\chi_j(R)$ 之积。
\end{definition}

\begin{theorem}[积分为零]\label{thm:b3:group-theory-applied:vanishing-integral}
对全空间的积分 $\int f_1f_2f_3\,\dd\tau$ 只有在
\omterm{def:b3:group-theory-applied:direct-product}{直积} $\Gamma_1\otimes\Gamma_2\otimes\Gamma_3$ 含有\omterm{def:b3:group-theory-applied:reducible}{全
对称表示}时才可能不为零。
\end{theorem}

\begin{proof}
\omterm{def:b3:point-groups:operation}{对称操作}只是对空间各点重新标记，故积分的值
在被积函数经任一 $R$ 变换后保持不变。取
变换后的被积函数在群上的平均：各\omterm{def:b3:point-groups:irrep}{不可约表示}所对应的
分量的平均为
$\frac1h\sum_R\hat R(\cdot)$，它是到\omterm{def:b3:group-theory-applied:reducible}{全对称
表示}上的投影（\cref{def:b3:group-theory-applied:projection}，取全部 $\chi = 1$）。
其余分量取平均后为零；若乘积不含
全对称部分，则积分为零。
\end{proof}

\begin{definition}[红外活性、拉曼活性]\label{def:b3:group-theory-applied:activity}
基频振动若能吸收红外光，则称其为\emph{红外活性}\index{红外活性}；
若其出现在拉曼光谱中，则称其为\emph{拉曼活性}\index{拉曼活性}
（见\cref{ch:b3:rovibrational-spectroscopy}）。
\end{definition}

\begin{corollary}[红外与拉曼选择定则]\label{cor:b3:group-theory-applied:ir-raman}
基频（由全对称的基能级跃迁到一个
对称性为 $\Gamma$ 的振动模式的量子）为\omterm{def:b3:group-theory-applied:activity}{红外活性}，仅当 $\Gamma$ 是
$x$、$y$ 或 $z$ 的表示；为\omterm{def:b3:group-theory-applied:activity}{拉曼活性}，仅当 $\Gamma$ 是
二次函数（$x^2$、$xy$、\dots）的表示。
\end{corollary}

\begin{proof}
吸收强度涉及积分 $\int\psi_0\,\mu_k\,\psi_1\dd\tau$，其中
偶极分量 $\mu_k$ 按 $x$、$y$、$z$ 变换；$\psi_0$ 为全对称，
$\psi_1$ 按 $\Gamma$ 变换。由定理，该积分非零仅当
$\Gamma\otimes\Gamma_{\mu_k}$ 含 $A_1$ 时成立，这只在 $\Gamma =
\Gamma_{\mu_k}$ 时成立（两个\omterm{def:b3:point-groups:irrep}{不可约表示}的乘积含\omterm{def:b3:group-theory-applied:reducible}{全对称表示}，
仅当二者相等（对实\omterm{def:b3:point-groups:representation}{特征标}）。拉曼
散射涉及极化率分量 $\alpha_{kl}$，其
变换方式与乘积 $kl$ 相同。
\end{proof}

\begin{proposition}[互斥规则]\label{prop:b3:group-theory-applied:mutual-exclusion}
在具有\omterm{def:b3:point-groups:inversion}{反演中心}的分子中，不存在同时既为\omterm{def:b3:group-theory-applied:activity}{红外活性}又为\omterm{def:b3:group-theory-applied:activity}{拉曼活性}的基频：
此即\emph{互斥规则}\index{互斥规则}。
\end{proposition}

\begin{proof}
有\omterm{def:b3:point-groups:inversion}{反演中心}时，每个\omterm{def:b3:point-groups:irrep}{不可约表示}均为 $g$ 或 $u$。
坐标 $x$、$y$、$z$ 在 $i$ 作用下改变符号（属 $u$）；二次
函数不改变符号（属 $g$）。一个模式仅当属 $u$ 时为\omterm{def:b3:group-theory-applied:activity}{红外活性}，
仅当属 $g$ 时为\omterm{def:b3:group-theory-applied:activity}{拉曼活性}。
\end{proof}

\begin{example}[二氧化碳与温室效应]\label{ex:b3:group-theory-applied:co2}
% ledger: vib:CO2.sym, vib:CO2.bend, vib:CO2.asym, ir:CO2.asym.int, ir:CO2.bend.int
\ce{CO2} 具有\omterm{def:b3:point-groups:inversion}{反演中心}。其对称伸缩振动（$\Sigma_g^+$，
\qty{1333}{cm^{-1}}）仅为\omterm{def:b3:group-theory-applied:activity}{拉曼活性}；反对称伸缩振动（$\Sigma_u^+$，
\qty{2349}{cm^{-1}}）和弯曲振动（$\Pi_u$，\qty{667}{cm^{-1}}）仅为\omterm{def:b3:group-theory-applied:activity}{红外活性}。
弯曲振动在 \qty{15}{\micro m} 附近吸收，接近地球热辐射的峰值，
这正是二氧化碳为温室气体的原因。\ce{N2} 与
\ce{O2} 各仅有一个 $\Sigma_g^+$ 振动，均无\omterm{def:b3:group-theory-applied:activity}{红外活性}。
\end{example}

\begin{omfigure}
\begin{tikzpicture}
\begin{axis}[omir, width=0.92\linewidth, height=4.4cm, xmin=400, xmax=2600,
  ymin=0, ymax=1.3, ytick=\empty, xlabel={波数 / cm$^{-1}$},
  ylabel={信号}, legend cell align=left,
  legend style={font=\small, at={(0.98,0.95)}, anchor=north east, draw=none}]
\addplot[omDef, thick] table[x=nu, y=ir] {figdata/out/bachelor-3-co2-spectra.dat};
\addlegendentry{红外吸收}
\addplot[omThm, thick, dashed] table[x=nu, y=raman] {figdata/out/bachelor-3-co2-spectra.dat};
\addlegendentry{拉曼散射}
\node[font=\small, anchor=south] at (axis cs:2349,1.02) {$\Sigma_u^+$};
\node[font=\small, anchor=south] at (axis cs:667,0.1) {$\Pi_u$};
\node[font=\small, anchor=south, omThm] at (axis cs:1333,1.02) {$\Sigma_g^+$};
\end{axis}
\end{tikzpicture}

{\small 气态 \ce{CO2} 的模拟光谱（谱带位置及相对
红外强度取自实测值；谱形仅为示意，不含
转动结构）。\omterm{prop:b3:group-theory-applied:mutual-exclusion}{互斥规则}：两谱中不出现同一谱带。}
\end{omfigure}

\begin{method}[预测振动光谱]\label{met:b3:group-theory-applied:vibrations}
\begin{enumerate}
  \item 求出\omterm{def:b3:point-groups:point-group}{点群}；由原位不动的原子计算 $\chi_{3N}$。
  \item 约化，扣除平动与转动：$\Gamma_{\mathrm{vib}}$。
  \item 由表的右侧各列，将各模式标为\omterm{def:b3:group-theory-applied:activity}{红外活性}
        （$x$、$y$、$z$）和/或\omterm{def:b3:group-theory-applied:activity}{拉曼活性}（二次函数）；两者皆非的模式
        不出现谱带。
  \item 数谱带：每个活性模式一条，简并模式只给出一条谱带。
\end{enumerate}
\end{method}

\begin{method}[计数 CO 伸缩振动]\label{met:b3:group-theory-applied:co-count}
在金属羰基化合物中，取 $n$ 个 C--O 键矢量为基（原位不动的原子计为
1，其余为零）：约化得 $\Gamma_{\mathrm{CO}}$，再应用
红外与拉曼选择定则。在
\qtyrange{1850}{2150}{cm^{-1}} 附近的强谱带数目可用于判定异构体。
\end{method}

\begin{example}[顺式与反式]\label{ex:b3:group-theory-applied:cis-trans}
顺式-\ce{ML4(CO)2}，$C_{2v}$：\omterm{def:b3:point-groups:representation}{特征标} $2, 0, 2, 0$，$\Gamma_{\mathrm{CO}} = A_1 + B_1$，
两条红外谱带。反式-\ce{ML4(CO)2}，$D_{4h}$：$\Gamma_{\mathrm{CO}} = A_{1g} + A_{2u}$，
一条红外谱带（$A_{2u}$）和一条拉曼谱带（$A_{1g}$）。红外区
仅见一条 CO 谱带，即表示为反式。
\end{example}

\section{电子跃迁的应用}

同一定理同样适用于电子跃迁：在态
$\Psi_1$ 与 $\Psi_2$ 之间的电子跃迁仅当 $\Gamma_1\otimes\Gamma_2$
包含 $x$、$y$ 或 $z$ 的表示时才被允许，且其光
沿该轴偏振。

\begin{example}[水与甲醛的跃迁]\label{ex:b3:group-theory-applied:electronic}
在水分子（$C_{2v}$，$xz$ 平面）中，电子由 $1b_2$（孤电子对）激发到
$4a_1$ 轨道，得到对称性为 $B_2\otimes A_1 = B_2$ 的激发态，即 $y$：
允许跃迁，偏振方向垂直于分子平面。在甲醛
（$C_{2v}$），$n\to\pi^*$ 跃迁由面内氧的孤电子对（$b_1$）
跃迁至 $\pi^*$ 轨道，该轨道由垂直于该平面的 $p$ 轨道组成
（$b_2$；分子位于 $xz$ 平面，C=O 沿 $z$）：$B_1\otimes B_2 = A_2$，
不属于 $x$、$y$、$z$ 中任一，即对称性禁阻，这正是其谱带
（见\cref{ch:b3:electronic-spectroscopy}）很弱的原因。
\end{example}

\begin{history}[维格纳、贝特与分子的对称性]
群论随尤金·维格纳关于原子光谱的工作（1927--31）进入量子力学，
以及汉斯·贝特关于原子谱项在晶体中分裂的论文（1929）。
E. 布赖特·威尔逊于
1934 年将其应用于分子振动；罗伯特·马利肯给出了至今仍在轨道与态中使用的标记。
\end{history}

\begin{inthelab}[红外与拉曼并排比较]
红外光谱仪测量样品吸收的光；拉曼光谱仪
用激光照射样品，分析其中
波长发生位移的弱散射光。对同一化合物同时记录两种谱，是
检验对称中心的经典方法：若无谱带重合，分子
很可能具有对称中心。拉曼激光器为 3B 类或 4 类，光束
被封闭，操作人员须佩戴与其波长相适配的护目镜。
\end{inthelab}

\section{练习}

\begin{exercise}[$\star$]\label{exo:b3:group-theory-applied:1}
约化 $C_{2v}$ 中\omterm{def:b3:point-groups:representation}{特征标}为 $(5, -1, 3, 1)$ 的表示，操作依次为 $(E, C_2,
\sigma_v(xz), \sigma_v'(yz))$。检验其维数。
\end{exercise}

\begin{exercise}[$\star$]\label{exo:b3:group-theory-applied:2}
\ce{SO2} 与水一样呈折线形。求 $\Gamma_{\mathrm{vib}}$，并指出哪些模式
具有红外和\omterm{def:b3:group-theory-applied:activity}{拉曼活性}。
\end{exercise}

\begin{exercise}[$\star$]\label{exo:b3:group-theory-applied:3}
在 $C_{2v}$ 中计算\omterm{def:b3:group-theory-applied:direct-product}{直积} $B_1\otimes B_2$、$A_2\otimes B_1$ 和
$E\otimes E$（在 $C_{3v}$ 中），并将后者约化。
\end{exercise}

\begin{exercise}[$\star$]\label{exo:b3:group-theory-applied:4}
% ledger: vib:CH4.nu1, vib:CH4.nu2, vib:CH4.nu3, vib:CH4.nu4
\ce{CH4} 的基频位于 2917（$a_1$）、1534（$e$）、3019 与 1306（均为 $t_2$）
\unit{cm^{-1}}。哪些出现在红外光谱中，
哪些出现在拉曼光谱中？
\end{exercise}

\begin{exercise}[$\star\star$]\label{exo:b3:group-theory-applied:5}
对 \ce{BF3}，由位置不变的原子推导 $\Gamma_{\mathrm{vib}} = A_1' + 2E' + A_2''$，
再统计红外和拉曼谱带数。哪一个模式只在两种谱
之一中出现，哪一个在两种谱中都出现？
\end{exercise}

\begin{exercise}[$\star\star$]\label{exo:b3:group-theory-applied:6}
\ce{SF6} 的红外和拉曼谱带各有多少条？哪一个模式
不出现于谱中？为什么两种谱中没有谱带同时出现？
\end{exercise}

\begin{exercise}[$\star\star$]\label{exo:b3:group-theory-applied:7}
对 $h_2$ 作用\omterm{def:b3:group-theory-applied:projection}{投影算符}，取 $E$ 表示（$C_{3v}$ 群），并
证明：将所得结果对 $2h_1 - h_2 - h_3$ 正交化后得 $h_2 - h_3$。
\end{exercise}

\begin{exercise}[$\star\star$]\label{exo:b3:group-theory-applied:8}
预测面式- 与经式-\ce{ML3(CO)3} 的红外 CO 谱带数，以及
\ce{M(CO)6} 的谱带数。
\end{exercise}

\begin{exercise}[$\star\star$]\label{exo:b3:group-theory-applied:9}
证明八面体配合物中配体的六个 $\sigma$ 轨道张成
$A_{1g} + E_g + T_{1u}$，并指出金属的哪些轨道可与各自结合。
\end{exercise}

\begin{exercise}[$\star\star\star$]\label{exo:b3:group-theory-applied:10}
正则表示满足 $\chi(E) = h$，且 $\chi(R) = 0$（$R \ne E$）。
用\omterm{thm:b3:group-theory-applied:reduction}{约化公式}证明：它包含每个\omterm{def:b3:point-groups:irrep}{不可约表示}
$\Gamma_i$ 恰 $d_i$ 次，并推出 $\sum_id_i^2 = h$。
\end{exercise}

\begin{exercise}[$\star\star\star$]\label{exo:b3:group-theory-applied:11}
乙炔 \ce{HC#CH} 为直线形。在 $D_{2h}$ 中求 $\Gamma_{\mathrm{vib}}$，
再将其重新归并为 $D_{\infty h}$ 的标记，并预测红外与拉曼光谱。
\end{exercise}

\begin{exercise}[$\star\star\star$]\label{exo:b3:group-theory-applied:12}
在水分子（$xz$ 平面）中，下列单电子激发中哪些给出
允许跃迁，并指出其偏振方向：$1b_2 \to 4a_1$，$3a_1 \to 4a_1$，
$1b_2 \to 2b_1$，$1b_1 \to 4a_1$？
\end{exercise}

\section{问题：顺式还是反式？由谱图判断羰基配合物}

\begin{problem}[{周末问题——点群与 CO 伸缩振动的
表示（四种异构体）、它们的红外与拉曼谱带，以及由实测
谱图鉴定产物}]
\label{pb:b3:group-theory-applied:1}
% ledger: ct:C2v, ct:C3v, ct:D4h
八面体金属 M 带有羰基配体和其他配体 L（每个 L
视为一个点）。考虑四种异构体：顺式-和
反式-\ce{ML4(CO)2}，以及面式- 与经式-\ce{ML3(CO)3}。假设 L 配体
除其位置所要求的之外，不降低对称性。合成
\ce{ML3(CO)3} 得到一种产物，其红外光谱显示三条强谱带，
位于 2048、1965 和 \qty{1938}{cm^{-1}}（题目数据）；其拉曼光谱
显示三条位置几乎相同的谱带。

\medskip\noindent\textbf{第一部分——\omterm{def:b3:point-groups:point-group}{点群}。}
\begin{enumerate}
  \item 在八面体上画出四种异构体。
  \item 给出顺式-\ce{ML4(CO)2} 的\omterm{def:b3:point-groups:point-group}{点群}。
  \item 反式-\ce{ML4(CO)2} 的\omterm{def:b3:point-groups:point-group}{点群}。
  \item 面式-\ce{ML3(CO)3} 的\omterm{def:b3:point-groups:point-group}{点群}。
  \item 经式-\ce{ML3(CO)3} 的\omterm{def:b3:point-groups:point-group}{点群}。
  \item 四种异构体中哪些具有\omterm{def:b3:point-groups:inversion}{反演中心}？
\end{enumerate}

\medskip\noindent\textbf{第二部分——CO 伸缩振动。}
\begin{enumerate}[resume]
  \item 说明 C--O 键矢量为何构成这样一种表示：
        其中每个保持原位的 CO 对操作贡献 1。
  \item 求顺式异构体的 $\Gamma_{\mathrm{CO}}$。
  \item 反式异构体。
  \item 面式异构体。
  \item 经式异构体。
  \item 将每个维数与 CO 配体的数目相核对。
\end{enumerate}

\medskip\noindent\textbf{第三部分——活性。}
\begin{enumerate}[resume]
  \item 统计每种异构体的\omterm{def:b3:group-theory-applied:activity}{红外活性} CO 伸缩振动数。
  \item 统计\omterm{def:b3:group-theory-applied:activity}{拉曼活性}的数目。
  \item 哪一种异构体符合\omterm{prop:b3:group-theory-applied:mutual-exclusion}{互斥规则}？
  \item 在反式异构体中描述\omterm{def:b3:group-theory-applied:activity}{红外活性}模式：两个 CO
        伸缩是同相还是反相？
  \item 在面式异构体中，三个 CO 为何只给出两条谱带？
  \item 红外谱带数为 1 说明什么？
\end{enumerate}

\medskip\noindent\textbf{第四部分——产物。}
\begin{enumerate}[resume]
  \item 产物的红外光谱表明是哪一种异构体？
  \item 拉曼光谱是否与之相符？
  \item 最高的谱带为 CO 基团的同相伸缩振动，在该异构体中它属于
        哪一个\omterm{def:b3:point-groups:irrep}{不可约表示}？
  \item 该谱图能否来自两种异构体的混合物？需要
        何种进一步测量才能判定？
  \item 给出结论：经式异构体\omterm{def:b3:group-theory-applied:activity}{红外活性} CO 伸缩振动的数目，
        与面式异构体的数目相比较。
\end{enumerate}
\end{problem}
